\section{Dark-matter couplings to a target}\label{sec:8.1}

A dark photon couples to a solid through its charge alone. In the language of \cref{sec:4.2} it needs one target operator per species, the electron density $n_e(\vect{q})$ and the density $n_\kappa(\vect{q})$ of each kind of nucleus, with potentials of opposite sign. A dark-matter particle with a magnetic dipole moment does more: it feels the spins of the electrons, and, since a charge moving past a magnetic moment feels a force, their velocities as well~\cite[Table~II]{trickle2022effective}. Particle physics offers many such models, and at first sight each would need its own calculation of the target. None does. A slow dark-matter particle sees an electron or a nucleus through three things only: the momentum it hands over, the velocity of the two particles relative to each other, and their spins. As a result every model reduces to the same short list of target operators per species: the density, the spin density, the current, the spin current and a kinetic-energy density. What remains of the model is a tensor on the dark-matter side that weights the response matrix of these operators.

\subsection{The non-relativistic operators}

Take a dark-matter particle of mass $\mchi$ and spin operator $\vect{S}_\chi$ that scatters on a free electron, with spin $\vect{S}_e = \boldsymbol{\sigma}/2$. The dark matter goes from $\vect{p}$ to $\vect{p}' = \vect{p} - \vect{q}$ and the electron from $\vect{k}$ to $\vect{k}' = \vect{k} + \vect{q}$, with $\vect{q}$ the momentum transfer of \cref{sec:1.3}. Momentum conservation leaves three of the four momenta independent, and Galilean invariance, the freedom to add a common velocity to both particles, allows an interaction to depend on them only through $\vect{q}$ and one velocity difference~\cite[Sect.~2.2]{fitzpatrick2013effective}. The difference that is Hermitian, symmetric between the incoming and outgoing particles, is
\begin{equation}
  \vect{v}^\perp = \bar{\vect{v}}_\chi - \frac{\vect{k} + \vect{k}'}{2m_e}, \qquad
  \bar{\vect{v}}_\chi \equiv \frac{\vect{p} + \vect{p}'}{2\mchi} = \vect{v} - \frac{\vect{q}}{2\mchi} ,
  \label{eq:8.1-vperp}
\end{equation}
the mean velocity of the dark matter minus that of the electron, with $\vect{v} = \vect{p}/\mchi$ the incoming velocity. For an electron initially at rest, $\vect{v}^\perp = \vect{v} - \vect{q}/2\mu_{\chi e}$, and energy conservation in the form $\vect{q}\cdot\vect{v} = q^2/2\mu_{\chi e}$ (\cref{sec:1.3}) makes $\vect{q}\cdot\vect{v}^\perp = 0$. This transversality is the origin of the name~\cite[Eq.~(14)]{fitzpatrick2013effective}, and \cref{app:D.24} shows that it holds for any free target. The momentum enters as $\ii\vect{q}$. Hermitian conjugation exchanges the incoming and outgoing particles, which reverses $\vect{q}$ and leaves $\vect{v}^\perp$ unchanged, so that $\ii\vect{q}$, $\vect{v}^\perp$ and the spins are all Hermitian.

An interaction is then a scalar built from $\ii\vect{q}$, $\vect{v}^\perp$, $\vect{S}_\chi$ and $\vect{S}_e$. Powers of $q^2$ can be absorbed into a momentum-dependent coefficient, as the mediator propagator of \cref{eq:1.4-Vq} already is, and products of more than one spin of the same particle reduce to at most one by the algebra of the Pauli matrices~\cite[Sect.~2.2]{fitzpatrick2013effective}. What remains is a finite list of operators, which Fitzpatrick and collaborators wrote down and Anand, Fitzpatrick and Haxton completed~\cites[Eqs.~(19)--(24) and (45)]{fitzpatrick2013effective}[Eqs.~(12) and (13)]{anand2014model}. They are listed in \cref{tab:8-operators}, which writes $\vect{q}/m_e$ in place of $\vect{q}$, so that each operator is dimensionless. The list is complete for dark matter of spin $\frac12$, whose spin products the Pauli algebra reduces. For spin 1 a product of two dark-matter spins leaves a symmetric tensor built from the polarization vectors, and two more operators appear, $\vect{q}$ contracted through that tensor with $\vect{v}^\perp$ or with $\vect{S}_e$~\cite[Eq.~(18)]{dent2015general}; for spin 0 only the operators free of $\vect{S}_\chi$ remain. The target side does not change: the new operators are linear in the electron's velocity or spin and reduce to the same rows. The potential between the dark matter and the electron is a sum
\begin{equation}
  \tilde V(\vect{q}) = \sum_i c_i^{(e)}(q)\,\mathcal{O}_i^{(e)} ,
  \label{eq:8.1-eft}
\end{equation}
an operator in the spin spaces of both particles, whose matrix element between plane waves is the transition amplitude in the first Born approximation (\cref{sec:2.2}). The coefficients $c_i^{(e)}(q)$ carry the dimension $\mathrm{eV}^{-2}$ of $\tilde V$ and the momentum dependence of the mediator. For the Yukawa potential of \cref{sec:1.4}, $c_1^{(e)}(q) = \tilde V(q) = y_\chi y_e/(q^2 + m_V^2)$, and every other coefficient vanishes. Papers that normalize their states relativistically quote $4\mchi m_e\,c_i^{(e)}$ instead, the factor of \cref{eq:A.9-M}. Refs.~\cite{fitzpatrick2013effective,anand2014model} define $\vect{q}$ as the momentum the dark matter gains, the opposite of the $\vect{q}$ of this book, so a coefficient quoted there changes sign for the operators odd in $\vect{q}$, $i = 3, 5, 9, 10, 11, 13, 14$; $\vect{v}^\perp$ is the same in both conventions.

Which operators a model produces follows from the non-relativistic limit of its couplings. The limits of the four Dirac bilinears are derived in \cref{app:C.15} and, between plane waves, in \cref{app:D.24}. The scalar $\bar\psi\psi$ and the time component of the vector $\bar\psi\gamma^\mu\psi$ both reduce to the density, $1$. The spatial vector reduces to the velocity plus a spin term, $(\vect{k} + \vect{k}')/2m_e - \ii\vect{q}\times\vect{S}_e/m_e$, the spatial axial vector $\bar\psi\gamma^\mu\gamma^5\psi$ to twice the spin, $2\vect{S}_e$, and the pseudoscalar $\bar\psi\ii\gamma^5\psi$ to $-\ii\vect{q}\cdot\vect{S}_e/m_e$. Multiplying a dark-matter bilinear by an electron bilinear and by the mediator propagator gives the coefficients. A vector mediator coupled to both vector currents leads to $\mathcal{O}_1$, and an axial-vector mediator coupled to both axial currents to $\mathcal{O}_4$, the standard spin-independent and spin-dependent interactions. A scalar mediator with a pseudoscalar coupling to the electron produces $\mathcal{O}_{10}$, and with pseudoscalar couplings on both sides $\mathcal{O}_6$. Dark matter with an electric dipole moment that couples through a photon-like mediator produces $\mathcal{O}_{11}$; a magnetic dipole moment, $\mathcal{O}_1$, $\mathcal{O}_4$, $\mathcal{O}_5$ and $\mathcal{O}_6$; and an anapole moment, $\mathcal{O}_8$ and $\mathcal{O}_9$~\cite[Table~II]{trickle2022effective}. The dipole and anapole couplings bring extra powers of $q$, in the operator or in the coefficient, and those coupled through a light mediator are strongly weighted toward small $q$. Operator $\mathcal{O}_2 = (\vect{v}^\perp)^2$ arises in none of these models at leading order~\cite[Sect.~2.1]{anand2014model}. It appears as a correction, of relative size $v^2$, to $\mathcal{O}_1$, and many lists omit it~\cite[Table~I]{catena2021crystal}; we keep it, because its reduction needs a row that the scalar absorption of \cref{sec:8.1-absorption} needs anyway.

\subsection{Target operators per species}

Inside a crystal the electron that the dark matter strikes is one of $10^{23}$, bound by the nuclei and by the other electrons, and its bare momentum $\vect{k}$ is not conserved; only its crystal momentum is (\cref{chap:3}). What survives is the structure of \cref{eq:8.1-eft}. Each operator acts on the target particle through three ingredients, the phase $\ee^{\ii\vect{q}\cdot\vect{r}}$ that adds the momentum $\vect{q}$, the velocity $(\vect{k} + \vect{k}')/2m_e$ inside $\vect{v}^\perp$, and the spin $\vect{S}_e$, and each has a one-body operator whose matrix elements between plane waves reproduce it. The phase is $\ee^{\ii\vect{q}\cdot\vect{r}}$ itself. The velocity is the symmetrized product $\tfrac12\{\vect{p}/m_e, \ee^{\ii\vect{q}\cdot\vect{r}}\}$, whose element from $\vect{k}$ to $\vect{k}'$ is $(\vect{k} + \vect{k}')/2m_e$, since $\vect{p}$ acting to the left gives $\vect{k}'$ and to the right gives $\vect{k}$. The square of the velocity is the doubly symmetrized $\{\vect{p}, \{\vect{p}, \ee^{\ii\vect{q}\cdot\vect{r}}\}\}/4m_e^2$. Summing over the particles, as in \cref{eq:2.3-onebody}, leads to five operators. Let $s$ be a species, either the electrons ($s = e$) or the nuclei of species $\kappa$ ($s = \kappa$), with mass $m_s$. Its particles $j$ have positions $\vect{r}_j$, momenta $\vect{p}_j$ and spins $\vect{S}_j$, and
\begin{keyresult}[Target operators per species]
\begin{equation}
  \begin{gathered}
  \begin{aligned}
  n_s(\vect{q}) &= \sum_{j\in s}\ee^{\ii\vect{q}\cdot\vect{r}_j}, &
  \vect{s}_s(\vect{q}) &= \sum_{j\in s}\ee^{\ii\vect{q}\cdot\vect{r}_j}\,\vect{S}_j, \\
  \vect{j}_s(\vect{q}) &= \sum_{j\in s}\frac12\Bigl\{\frac{\vect{p}_j}{m_s},\,\ee^{\ii\vect{q}\cdot\vect{r}_j}\Bigr\}, &
  \mathcal{J}^{\alpha\beta}_s(\vect{q}) &= \sum_{j\in s}\frac12\Bigl\{\frac{p^\alpha_j}{m_s},\,\ee^{\ii\vect{q}\cdot\vect{r}_j}\Bigr\}\,S^\beta_j,
  \end{aligned}\\
  \mathcal{T}_s(\vect{q}) = \sum_{j\in s}\frac{1}{8m_s}\bigl\{\vect{p}_j, \{\vect{p}_j, \ee^{\ii\vect{q}\cdot\vect{r}_j}\}\bigr\}
  = \sum_{j\in s}\frac{\vect{p}_j\cdot\ee^{\ii\vect{q}\cdot\vect{r}_j}\,\vect{p}_j}{2m_s} + \frac{q^2}{8m_s}\,n_s(\vect{q}) .
  \end{gathered}
  \label{eq:8.1-rows}
\end{equation}
These are the density, the spin density, the paramagnetic current, the spin current (velocity index $\alpha$, spin index $\beta$) and the kinetic-energy density. Their matrix elements from a plane wave $\vect{k}$ to $\vect{k} + \vect{q}$ are $1$, $\vect{S}$, $(\vect{k} + \vect{k}')/2m_s$, $(k^\alpha + k'^\alpha)S^\beta/2m_s$ and $(\vect{k} + \vect{k}')^2/8m_s$.
\end{keyresult}
\noindent For the electrons they are old acquaintances. The density $n_e(\vect{q})$ is \cref{eq:2.3-nq}, the spin density $\vect{s}_e(\vect{q})$ is \cref{eq:7.1-se}, and the current $\vect{j}_e(\vect{q})$ is that of the continuity equation \cref{eq:4.3-continuity}. The spin current and the kinetic-energy density are new. At $\vect{q} = 0$ the kinetic-energy density is the kinetic energy $\sum_j p_j^2/2m_s$ of the species, and at finite $\vect{q}$ it is the kinetic energy of the mean of the incoming and outgoing momenta, which is why the term $q^2/8m_s$ appears beside the familiar $\vect{p}\cdot\ee^{\ii\vect{q}\cdot\vect{r}}\vect{p}/2m_s$. For nuclei the spin $\vect{S}_j$ is the nuclear spin, and \cref{sec:8.1-nuclei} turns couplings to nucleons into couplings to these rows.

With the rows in hand the reduction is mechanical. Substitute $\vect{v}^\perp = \bar{\vect{v}}_\chi - (\vect{k} + \vect{k}')/2m_e$ in each operator, expand, and replace every monomial in the electron's velocity and spin by its row: $1 \to n_e$, $\vect{S}_e \to \vect{s}_e$, the velocity $\to \vect{j}_e$, the velocity times the spin $\to \mathcal{J}_e$, and the square of the velocity $\to (2/m_e)\mathcal{T}_e$. The dark-matter velocity $\bar{\vect{v}}_\chi$ and spin $\vect{S}_\chi$ are not target operators; they stay in the coefficients. For instance,
\begin{equation}
  \mathcal{O}_8 = \vect{S}_\chi\cdot\vect{v}^\perp \;\to\; (\vect{S}_\chi\cdot\bar{\vect{v}}_\chi)\,n_e(\vect{q}) - \vect{S}_\chi\cdot\vect{j}_e(\vect{q}) ,
  \qquad
  \mathcal{O}_7 = \vect{S}_e\cdot\vect{v}^\perp \;\to\; \bar{\vect{v}}_\chi\cdot\vect{s}_e(\vect{q}) - \mathcal{J}^{\alpha\alpha}_e(\vect{q}) ,
  \label{eq:8.1-O7O8}
\end{equation}
with a sum over the repeated Cartesian index. \Cref{tab:8-operators} lists all fifteen reductions, and \cref{app:D.24} derives them. Every operator except $\mathcal{O}_2$ is at most linear in the electron's velocity and in its spin, and needs at most the four rows $n_e$, $\vect{s}_e$, $\vect{j}_e$ and $\mathcal{J}_e$~\cite[Table~III]{trickle2022effective}; $\mathcal{O}_2$ adds $\mathcal{T}_e$.

\begin{table}[tbp]
  \caption{Non-relativistic operators of dark-matter scattering and absorption, and the target operators each needs. Upper block: the operators $\mathcal{O}_i$ of Refs.~\cite[Eqs.~(19)--(24) and (45)]{fitzpatrick2013effective} and~\cite[Eqs.~(12) and (13)]{anand2014model}, written for an electron with $\vect{q}$ the momentum the dark matter loses; $\vect{v}^\perp$ is \cref{eq:8.1-vperp}, $\bar{\vect{v}} \equiv \bar{\vect{v}}_\chi = \vect{v} - \vect{q}/2\mchi$, and $m = m_e$. The rows are those of \cref{eq:8.1-rows} for $s = e$, all at the same $\vect{q}$; repeated Cartesian indices are summed, and $\epsilon_{\alpha\beta\gamma}$ is the Levi-Civita symbol. For nucleons, replace $e$ by $N$ and $m_e$ by the nucleon mass; the nuclear rows follow from \cref{eq:8.1-nuclear}. Lower block: absorption of a boson $\phi$ of mass $m_\phi$ at $\vect{q} = \vect{p}_\phi$, and fermionic absorption with an outgoing fermion $\tilde\nu$ (\cref{sec:8.1-absorption}). Last column: what each row excites in an insulator. Through the nuclei the density row also excites phonons, and in an ordered magnet the spin rows excite magnons (\cref{sec:7.1}).}
  \label{tab:8-operators}
  \centering\small
  \renewcommand{\arraystretch}{1.25}
  \begin{tabular}{@{}l>{\raggedright\arraybackslash}p{0.25\linewidth}>{\raggedright\arraybackslash}p{0.39\linewidth}>{\raggedright\arraybackslash}p{0.15\linewidth}@{}}
    \toprule
    & Operator & Target operators & Excites \\
    \midrule
    $\mathcal{O}_1$ & $1$ & $n$ & pairs, plasmons \\
    $\mathcal{O}_2$ & $(\vect{v}^\perp)^2$ & $\bar v^2 n - 2\bar{\vect{v}}\cdot\vect{j} + (2/m)\mathcal{T}$ & pairs \\
    $\mathcal{O}_3$ & $\ii\vect{S}_e\cdot(\vect{q}/m\times\vect{v}^\perp)$ & $(\ii/m)\bigl[(\vect{q}\times\bar{\vect{v}})\cdot\vect{s} - \epsilon_{\alpha\beta\gamma}q^\beta\mathcal{J}^{\gamma\alpha}\bigr]$ & spin flips \\
    $\mathcal{O}_4$ & $\vect{S}_\chi\cdot\vect{S}_e$ & $\vect{S}_\chi\cdot\vect{s}$ & spin flips \\
    $\mathcal{O}_5$ & $\ii\vect{S}_\chi\cdot(\vect{q}/m\times\vect{v}^\perp)$ & $(\ii/m)\,\vect{S}_\chi\cdot\bigl[(\vect{q}\times\bar{\vect{v}})\,n - \vect{q}\times\vect{j}\bigr]$ & pairs \\
    $\mathcal{O}_6$ & $(\vect{S}_\chi\cdot\vect{q}/m)(\vect{S}_e\cdot\vect{q}/m)$ & $(\vect{S}_\chi\cdot\vect{q})(\vect{q}\cdot\vect{s})/m^2$ & spin flips \\
    $\mathcal{O}_7$ & $\vect{S}_e\cdot\vect{v}^\perp$ & $\bar{\vect{v}}\cdot\vect{s} - \mathcal{J}^{\alpha\alpha}$ & spin flips \\
    $\mathcal{O}_8$ & $\vect{S}_\chi\cdot\vect{v}^\perp$ & $(\vect{S}_\chi\cdot\bar{\vect{v}})\,n - \vect{S}_\chi\cdot\vect{j}$ & pairs \\
    $\mathcal{O}_9$ & $\ii\vect{S}_\chi\cdot(\vect{S}_e\times\vect{q}/m)$ & $(\ii/m)\,\vect{S}_\chi\cdot(\vect{s}\times\vect{q})$ & spin flips \\
    $\mathcal{O}_{10}$ & $\ii\vect{S}_e\cdot\vect{q}/m$ & $(\ii/m)\,\vect{q}\cdot\vect{s}$ & spin flips \\
    $\mathcal{O}_{11}$ & $\ii\vect{S}_\chi\cdot\vect{q}/m$ & $(\ii/m)(\vect{S}_\chi\cdot\vect{q})\,n$ & pairs \\
    $\mathcal{O}_{12}$ & $\vect{S}_\chi\cdot(\vect{S}_e\times\vect{v}^\perp)$ & $\vect{S}_\chi\cdot(\vect{s}\times\bar{\vect{v}}) - \epsilon_{\alpha\beta\gamma}S^\alpha_\chi\mathcal{J}^{\gamma\beta}$ & spin flips \\
    $\mathcal{O}_{13}$ & $\ii(\vect{S}_\chi\cdot\vect{v}^\perp)(\vect{S}_e\cdot\vect{q}/m)$ & $(\ii/m)\bigl[(\vect{S}_\chi\cdot\bar{\vect{v}})(\vect{q}\cdot\vect{s}) - S^\alpha_\chi q^\beta\mathcal{J}^{\alpha\beta}\bigr]$ & spin flips \\
    $\mathcal{O}_{14}$ & $\ii(\vect{S}_\chi\cdot\vect{q}/m)(\vect{S}_e\cdot\vect{v}^\perp)$ & $(\ii/m)(\vect{S}_\chi\cdot\vect{q})\bigl(\bar{\vect{v}}\cdot\vect{s} - \mathcal{J}^{\alpha\alpha}\bigr)$ & spin flips \\
    $\mathcal{O}_{15}$ & $-(\vect{S}_\chi\cdot\vect{q}/m)\bigl[(\vect{S}_e\times\vect{v}^\perp)\cdot\vect{q}/m\bigr]$ & $-(\vect{S}_\chi\cdot\vect{q})\bigl[\vect{q}\cdot(\vect{s}\times\bar{\vect{v}}) - \epsilon_{\alpha\beta\gamma}q^\alpha\mathcal{J}^{\gamma\beta}\bigr]/m^2$ & spin flips \\
    \midrule
    Dark photon & $-\zeta e\,A'_\mu J^\mu_{\mathrm{em}}$ & $-\vect{j}_e + \sum_\kappa Z_\kappa\vect{j}_\kappa$ & pairs; polar phonons \\
    ALP & $g_{aee}\,\phi\,\bar\psi\ii\gamma^5\psi$ & $\vect{s}_e$ ($\propto|\vect{p}_\phi|$); $\mathcal{J}^{\alpha\alpha}_e$ ($\propto m_\phi$) & pairs \\
    Scalar & $g_{\phi ee}\,\phi\,\bar\psi\psi$ & $n_e - \mathcal{T}_e/m_e$ & pairs \\
    Fermionic & $\Lambda_F^{-2}(\bar{\tilde\nu}\Gamma\chi)(\bar\psi\Gamma'\psi)$ & $n$, $\vect{j}$, $\vect{s}$, $\mathcal{J}^{\alpha\alpha}$ or $n - \mathcal{T}/m$, by $\Gamma'$ & pairs; nuclear recoils \\
    \bottomrule
  \end{tabular}
\end{table}

In a solid the electron moves too. The velocity-dependent operators split into a piece in which the dark-matter velocity multiplies a velocity-independent row and a piece that needs the target's own velocity. For a free electron at rest the second piece vanishes, but in a solid it does not, and it is often the larger: the valence electrons of silicon, taken as a free gas at their density, move at $v_{\mathrm{F}} = k_{\mathrm{F}}/m_e = 7.0\times10^{-3}$ (\cref{sec:4.5}), faster than any dark-matter particle in the laboratory. In the language of Ref.~\cite[Table~III]{trickle2022effective}, the target velocity turns a coupling to the charge or spin of the electron into a coupling to its orbital motion. The rows also overlap, in the longitudinal direction. The longitudinal current is not independent of the density: by the continuity equation \cref{eq:4.3-continuity}, $\langle f|\vect{q}\cdot\vect{j}_e|i\rangle = (E_f - E_i)\langle f|n_e|i\rangle$ between eigenstates of any Hamiltonian whose potentials are local. On the dark-matter side, $\vect{q}\cdot\bar{\vect{v}}_\chi = \vect{q}\cdot\vect{v} - q^2/2\mchi = \omega + \Delta_\chi$ by \cref{eq:1.3-omegaq}. In every matrix element that enters a rate, therefore,
\begin{equation}
  \vect{q}\cdot\vect{v}^\perp \;\to\; \Delta_\chi\,n_e(\vect{q}) ,
  \label{eq:8.1-transverse}
\end{equation}
which vanishes for elastic dark matter (\cref{app:D.24}). The transversality of $\vect{v}^\perp$ is thus not a property of free particles alone: as long as the target's potentials are local, only the components of $\bar{\vect{v}}_\chi$ and $\vect{j}_e$ transverse to $\vect{q}$ enter the rate of an elastic process. The same argument applies to the spin current, $q^\alpha\mathcal{J}^{\alpha\beta}_e \to \omega\,s^\beta_e$, provided the Hamiltonian conserves each electron's spin, which spin--orbit coupling and the hyperfine interaction violate.

Collecting the coefficients, the interaction between a dark-matter particle and the target, from dark-matter momentum $\vect{p}$ to $\vect{p} - \vect{q}$, takes the form assumed in \cref{sec:4.2},
\begin{equation}
  \langle\vect{p} - \vect{q}|H_{\mathrm{int}}|\vect{p}\rangle = \frac1\Vol\sum_a\tilde V_a(\vect{q}, \bar{\vect{v}}_\chi)\,O_a(\vect{q}) ,
  \label{eq:8.1-Hint}
\end{equation}
with the index $a$ running over the rows of every species. The only change from \cref{sec:4.2} is that the potentials $\tilde V_a$ now depend on the dark-matter velocity and are matrices in the spin space of the dark matter. For $\mathcal{O}_8$, for instance, $\tilde V_{n_e} = c_8^{(e)}\,\vect{S}_\chi\cdot\bar{\vect{v}}_\chi$ and $\tilde V_{j^\alpha_e} = -c_8^{(e)}S^\alpha_\chi$. An unpolarized beam averages over the initial spin of the dark matter and sums over its final spin, and the square of the amplitude $\Vol^{-1}\sum_b\langle f|O_b|i\rangle\tilde V_b$ becomes a trace over that spin space:
\begin{keyresult}[Dark-matter-side tensor]
\begin{equation}
  \mathfrak{D}_{ab}(\vect{q}, \vect{v}) = \frac{1}{2|\vect{S}_\chi| + 1}\,\Tr_\chi\bigl[\tilde V_a(\vect{q}, \bar{\vect{v}}_\chi)^\dagger\,\tilde V_b(\vect{q}, \bar{\vect{v}}_\chi)\bigr] ,
  \label{eq:8.1-D}
\end{equation}
with $|\vect{S}_\chi|$ the spin quantum number of the dark matter. The rate of \cref{eq:4.2-rate} holds with $\tilde V_a^*\tilde V_b \to \mathfrak{D}_{ab}$. The tensor is Hermitian and positive semidefinite, and for a coupling that does not involve $\vect{S}_\chi$ it is $\tilde V_a^*\tilde V_b$.
\end{keyresult}
\noindent Since $\Tr_\chi\vect{S}_\chi = 0$, a coefficient linear in $\vect{S}_\chi$ never interferes with one that is independent of it, and the operators split into two classes that add in the rate. For spin $\frac12$, $\Tr_\chi(S^\alpha_\chi S^\beta_\chi)/2 = \delta_{\alpha\beta}/4$.

Responses of the density are what \cref{chap:4} computed. The other rows need the same calculation with other vertices. For the spin density and the spin current, the vertex is that of the density or of the current multiplied by a Pauli matrix. For the current and the kinetic-energy density it needs the momentum operator, whose matrix elements change when the pseudopotential is nonlocal and must be corrected for it. In practice a calculation computes the independent-particle response matrix of all the electron rows at once and screens it with one Dyson equation (\cref{sec:8.4}); \cref{app:E.16} describes the tables it exports.

\subsection{Nuclei}\label{sec:8.1-nuclei}

A model specifies its couplings to protons and neutrons, $c_i^{(p)}$ and $c_i^{(n)}$, with operators $\mathcal{O}_i^{(N)}$ built from the nucleon spin $\vect{S}_N$, the nucleon mass $m_N$ and the nucleon velocity. The solid, however, contains nuclei. Nucleus $I$ of species $\kappa$ sits at $\vect{X}_I$, and its $A_\kappa$ nucleons sit at $\vect{X}_I + \vect{y}_j$, within a radius of a few fm. The small size of the nucleus and the small energy deposited reduce the nucleon operators to nuclear rows. Light dark matter transfers at most $2\mchi v_{\mathrm{max}} \approx \SI{5}{MeV}$, against $1/R \approx \SI{57}{MeV}$ for silicon (\cref{sec:1.4}), so $\vect{q}\cdot\vect{y}_j$ is small. And the energy it deposits on a nucleus, at most \SI{0.50}{keV} on silicon at $\mchi = \SI{1}{GeV}$ (\cref{sec:1.3}), lies far below the excitation energies of nuclei~\cite[Sect.~2.2]{fitzpatrick2013effective}, so the nucleus stays in its ground multiplet.
The small value of $\vect{q}\cdot\vect{y}_j$ turns the phase $\ee^{\ii\vect{q}\cdot(\vect{X}_I + \vect{y}_j)}$ into $\ee^{\ii\vect{q}\cdot\vect{X}_I}$ times the nuclear form factor $F_N(q)$ of \cref{eq:2.2-FN}, and the nucleon couplings add coherently, as in \cref{sec:1.4}. The nucleon velocity is the velocity of the nucleus, $\vect{P}_I/M_\kappa$, plus an internal velocity; the internal part averages to zero in the ground state and contributes at higher order in $\vect{q}\cdot\vect{y}_j$~\cite[Eqs.~(26) and (27)]{fitzpatrick2013effective}, and we drop it. The sum of the nucleon spins, weighted by their couplings, is a vector operator, and since the nucleus stays in its ground multiplet, the Wigner--Eckart theorem makes it proportional to the nuclear spin $\vect{J}^{\mathrm{nuc}}_I$ within the ground multiplet of spin $J_{\mathrm{nuc}}$. The constant follows from the state with $M = J_{\mathrm{nuc}}$, in which the proton and neutron spins have the expectation values $\langle S_p\rangle$ and $\langle S_n\rangle$ along the nuclear spin.
Together these give, for nucleus species $\kappa$ (and, for spin couplings, a single isotope),
\begin{equation}
  c_i^{(\kappa)}(q) = \bigl[Z_\kappa c_i^{(p)} + (A_\kappa - Z_\kappa)\,c_i^{(n)}\bigr]F_N(q)
  \quad\text{or}\quad
  c_i^{(\kappa)}(q) = \frac{c_i^{(p)}\langle S_p\rangle + c_i^{(n)}\langle S_n\rangle}{J_{\mathrm{nuc}}} ,
  \label{eq:8.1-nuclear}
\end{equation}
the first for operators without $\vect{S}_N$ and the second, with $\vect{S}_N \to \vect{J}^{\mathrm{nuc}}_I$, for those linear in it. The nuclear operators then reduce as in \cref{tab:8-operators}, to the rows of \cref{eq:8.1-rows} with $s = \kappa$, $m_s = M_\kappa$ and $\vect{S}_j = \vect{J}^{\mathrm{nuc}}_I$; in particular the nuclear spin density is $\vect{s}_\kappa(\vect{q}) = \sum_{I\in\kappa}\ee^{\ii\vect{q}\cdot\vect{X}_I}\vect{J}^{\mathrm{nuc}}_I$. The operators keep $\vect{q}/m_N$ with the nucleon mass, since that is how the coefficients are defined.

The spin coupling is far weaker than the coherent one. Silicon has one isotope with a nuclear spin, $^{29}$Si, with $J_{\mathrm{nuc}} = \frac12$~\cite[Fig.~5]{klos2013large} and an abundance of \SI{4.685}{\percent} in natural silicon~\cite{meija2016isotopic}. Shell-model calculations give $\langle S_n\rangle = 0.156$ and $\langle S_p\rangle = 0.016$~\cite[Table~III]{klos2013large}. For a coupling to neutron spins alone, \cref{eq:8.1-nuclear} gives $c_4^{(\kappa)} = 0.312\,c_4^{(n)}$. Averaged over the initial orientations of the nucleus and summed over the final ones, the squared spin operator is $J_{\mathrm{nuc}}(J_{\mathrm{nuc}} + 1) = \frac34$. The strength is then $(0.312)^2 \times \frac34 = 0.073$ in units of $|c_4^{(n)}|^2$, against $\frac34$ for a free neutron. A $^{29}$Si nucleus thus responds with \SI{10}{\percent} of the strength of one neutron, whereas its spin-independent coupling is $A^2 = 841$ times that of one nucleon, and in natural silicon one nucleus in 21 carries the spin. A sample of unpolarized nuclear spins has no coherent spin response, because the spins of two different nuclei are uncorrelated; only the incoherent part of $S_{\vect{s}_\kappa\vect{s}_\kappa}$ survives (\cref{sec:11.6}).

\begin{workedexample}{A heavy dark photon reduces to the density rows}\label{we:8-darkphoton}
Take the thermal benchmark of \cref{we:1-darkphoton}: $\mchi = \SI{10}{MeV}$, $m_V = \SI{30}{MeV}$, $y_\chi = 0.3$ and $\zeta = 10^{-4}$. The dark photon couples to the dark-matter current with strength $y_\chi$ and to $J^\mu_{\mathrm{em}}$ with strength $\zeta e$ (\cref{sec:1.4}). Exchanging it between a dark-matter particle and a target fermion with the currents $J^\mu_\chi$ and $J^\mu$ gives, for $\omega \ll q$, the potential $y_\chi y\,(J^0_\chi J^0 - \vect{J}_\chi\cdot\vect{J})/(q^2 + m_V^2)$, with $y$ the coupling of the target fermion (\cref{app:D.24}). The time components reduce to 1 on both sides, and the leading term is $\mathcal{O}_1$ with
\begin{equation*}
  c_1^{(e)} = \frac{y_\chi y_e}{q^2 + m_V^2} = \tilde V(q), \qquad c_1^{(p)} = -c_1^{(e)}, \qquad c_1^{(n)} = 0 ,
\end{equation*}
since $y_e = -\zeta e$ and a proton has $+\zeta e$. By \cref{eq:8.1-nuclear}, $c_1^{(\kappa)} = -Z_\kappa\tilde V F_N(q)$, and with $F_N = 1$ to better than \SI{0.1}{\percent} at these momenta the interaction is $\Vol^{-1}\tilde V(q)\bigl[n_e(\vect{q}) - \sum_\kappa Z_\kappa n_\kappa(\vect{q})\bigr] = -\Vol^{-1}\tilde V(q)\,n_Z(\vect{q})$, with $n_Z$ the charge-number density of \cref{sec:2.3}. The tensor is $\mathfrak{D}_{ab} = \tilde V_a^*\tilde V_b$ over the rows $n_e$ and $n_\kappa$, and the rate is $|\tilde V|^2$ times the structure factor of the total charge. Its electron block is the $S$ of \cref{eq:1.1-rate}; the nuclear block and the cross terms, of opposite sign, matter where the nuclei move, at phonon energies (\cref{sec:8.4}).

What is dropped? At $m_V = \SI{30}{MeV}$ the mediator is heavy at every momentum transfer, $q \le 2\mchi v_{\mathrm{max}} = \SI{53}{keV}$, and the propagator is constant to $(q/m_V)^2 \le 3\times10^{-6}$. The largest correction to the density vertex is the factor $1 - q^2/8m_e^2$ of \cref{eq:C.15-vertex}, which is $7\times10^{-6}$ at $q_0$ and $1.3\times10^{-3}$ at \SI{53}{keV}. The spatial currents, $\vect{J}_\chi\cdot\vect{J}$, are products of velocities and spin terms. Their spin-independent piece, $\bar{\vect{v}}_\chi\cdot(\vect{k} + \vect{k}')/2m_e$, is at most $v_{\mathrm{max}}v_{\mathrm{F}} = 1.9\times10^{-5}$ for valence electrons. The pieces linear in one spin do not interfere with the density coupling in a nonmagnetic target and enter the rate squared, and the product of the two spin terms is at most $q^2/(4\mchi m_e) = 1.4\times10^{-4}$. To a part in $10^3$, therefore, a heavy dark photon couples to the density rows and nothing else.
\end{workedexample}

\subsection{The velocity-squared row}\label{sec:8.1-v2}

The operator $\mathcal{O}_2 = (\vect{v}^\perp)^2$ is the only one that is quadratic in a velocity. Expanding the square,
\begin{equation}
  (\vect{v}^\perp)^2 = \bar v_\chi^2 - 2\bar{\vect{v}}_\chi\cdot\frac{\vect{k} + \vect{k}'}{2m_e} + \frac{(\vect{k} + \vect{k}')^2}{4m_e^2}
  \;\to\; \bar v_\chi^2\,n_e(\vect{q}) - 2\bar{\vect{v}}_\chi\cdot\vect{j}_e(\vect{q}) + \frac{2}{m_e}\,\mathcal{T}_e(\vect{q}) .
  \label{eq:8.1-O2}
\end{equation}
The first two terms use the density and current rows. The third needs the square of the electron's velocity, and since $\mathcal{O}_2$ is a scalar, only the trace of the velocity tensor $v^\alpha v^\beta$ appears, which is $2/m_e$ times the kinetic-energy density of \cref{eq:8.1-rows}. One row per species therefore suffices.

How large is each term? A naive estimate sets $\vect{v}^\perp \sim \vect{v} \sim 10^{-3}$ and makes $\mathcal{O}_2$ a correction of about $10^{-6}$ to $\mathcal{O}_1$. That estimate, however, treats the electron as being at rest, and \cref{eq:8.1-O2} shows what it misses. By \cref{eq:8.1-transverse}, the components along $\vect{q}$ of the dark-matter and electron velocities cancel for elastic scattering, and what remains of $\vect{v}^\perp$ is the difference of their transverse parts. The dark matter's transverse velocity obeys $v_{\mathrm{T}}^2 = v^2 - \vmin^2 \le v_{\mathrm{max}}^2$, since its longitudinal part is $\vmin$ at $\Delta_\chi = 0$. The electron's is set by the motion of the electrons in the solid, about $v_{\mathrm{F}}$ for valence states and larger for the core.

\begin{workedexample}{The velocity-squared operator in silicon}\label{we:8-O2}
We model the valence electrons of silicon as the free electron gas of \cref{sec:4.5}, at $\nd_e = \SI{2.00e23}{cm^{-3}}$, with $k_{\mathrm{F}} = \SI{3.57}{keV}$ and $v_{\mathrm{F}} = 6.99\times10^{-3}$, at $T = 0$. Take dark matter of mass \SI{100}{MeV} moving at $v = 2.0\times10^{-3}$ that transfers $q = q_0 = \SI{3.73}{keV}$ and deposits $\omega = \SI{5}{eV}$. A pair takes an electron from $\vect{k}$ inside the Fermi sphere to $\vect{k} + \vect{q}$ outside it, with $\omega = \vect{q}\cdot(2\vect{k} + \vect{q})/2m_e$. This fixes the component of $\vect{k}$ along $\vect{q}$ at $m_e\omega/q - q/2 = \SI{-1.18}{keV}$. The two Pauli conditions then confine the transverse momentum $k_{\mathrm{T}}$ to an annulus, $k_{\mathrm{F}}^2 - (\SI{2.55}{keV})^2 < k_{\mathrm{T}}^2 < k_{\mathrm{F}}^2 - (\SI{1.18}{keV})^2$, with \SI{2.55}{keV} the longitudinal momentum after the transfer, that is $\SI{2.50}{keV} < k_{\mathrm{T}} < \SI{3.37}{keV}$, which every pair of this $(\vect{q}, \omega)$ fills uniformly in area. The electron's transverse velocity, $(\vect{k} + \vect{k}')_{\mathrm{T}}/2m_e = \vect{k}_{\mathrm{T}}/m_e$, has $\langle v_{e\mathrm{T}}^2\rangle = 3.37\times10^{-5}$ and $\langle v_{e\mathrm{T}}^4\rangle = 1.17\times10^{-9}$ over the annulus.

On the dark-matter side, $\vmin = \omega/q + q/2\mchi = 1.36\times10^{-3}$ and $v_{\mathrm{T}}^2 = v^2 - \vmin^2 = 2.15\times10^{-6}$, smaller than the electron's by a factor of 16; at $v = v_{\mathrm{max}} = 2.650\times10^{-3}$ the factor is still 6.5. Per pair, $(\vect{v}^\perp)^2 = |\vect{v}_{\mathrm{T}} - \vect{v}_{e\mathrm{T}}|^2$, and the azimuth of $\vect{v}_{e\mathrm{T}}$ about $\vect{q}$ is uniform. Averaging its square, with $A = v_{\mathrm{T}}^2$ and $B = v_{e\mathrm{T}}^2$, gives $(A + B)^2 + 2AB$, and over the annulus
\begin{equation*}
  \bigl\langle(\vect{v}^\perp)^4\bigr\rangle = A^2 + 4A\langle B\rangle + \langle B^2\rangle
  = (0.0046 + 0.290 + 1.167)\times10^{-9} = 1.46\times10^{-9} .
\end{equation*}
The term that the naive estimate keeps, $A^2$, is \SI{0.3}{\percent} of the total. The fourth power of the electron's transverse velocity, $\langle B^2\rangle$, carries \SI{80}{\percent}, and the cross term $4A\langle B\rangle$, in which the two transverse velocities meet, the remaining \SI{20}{\percent}. These are shares of velocities, not of rows: every row of \cref{eq:8.1-O2} also carries longitudinal velocity, which cancels in their sum. At equal coefficients the rate of $\mathcal{O}_2$ is $1.5\times10^{-9}$ times that of $\mathcal{O}_1$ at this $(\vect{q}, \omega)$, about 90 times the naive $v^4 = 1.6\times10^{-11}$.
\end{workedexample}

The electron gas understates the effect in a real crystal. The valence wavefunctions of silicon have Fourier components at the reciprocal-lattice vectors (\cref{sec:3.2}), so the electron velocities spread beyond $v_{\mathrm{F}}$, and the $1s$ electrons, whose absorption edge lies at \SI{1.84}{keV} (\cref{sec:4.3}), move at $Z\alpha \approx 0.10$, where $(\vect{v}^\perp)^2 \approx 10^{-2}$ (\cref{app:C.15}). A velocity-dependent coupling therefore weights the core and the high-momentum tails of the valence states far more than the density coupling does. In practice, it is the target term $\mathcal{T}_e$ that a calculation must get right.

\subsection{Absorption operators}\label{sec:8.1-absorption}

A bosonic dark-matter field $\phi$ of mass $m_\phi$ oscillates as $\ee^{-\ii m_\phi t + \ii\vect{p}_\phi\cdot\vect{x}}$, and absorbing one quantum deposits $\omega = m_\phi$ and $\vect{q} = \vect{p}_\phi$ (\cref{eq:1.4-absorption}). Its coupling to the target is a bilinear of the target fields multiplied by $\phi$ or a derivative of $\phi$, and the spatial integral over $\ee^{\ii\vect{p}_\phi\cdot\vect{x}}$ turns each bilinear into a row of \cref{eq:8.1-rows} at $\vect{q} = \vect{p}_\phi$. The rows are the same as for scattering. What differs is the point $(\vect{q}, \omega)$ at which their response is needed, $|\vect{p}_\phi| \approx 10^{-3}\omega$ instead of $\omega \le qv$. The three bosons of \cref{sec:1.4} need different rows (\cref{app:D.24}).

The dark photon couples to $J^\mu_{\mathrm{em}}$ like a photon, with $e \to \zeta e$. For a dark photon $A'_\mu$ of mass $m_V = m_\phi$ at rest, the time component is smaller than the spatial ones by $|\vect{p}_\phi|/m_\phi$, by the Lorenz condition, and the coupling is to the spatial current of all charges, $-\vect{j}_e(\vect{p}_\phi) + \sum_\kappa Z_\kappa\vect{j}_\kappa(\vect{p}_\phi)$. At $\vect{p}_\phi \to 0$ its response is the optical conductivity (\cref{sec:8.2-absorption}). The nuclear currents carry the absorption into polar optical phonons below the electronic gap.

The axion-like particle couples through \cref{eq:1.4-alp}, in which we write $\phi$, $\nabla\phi$ and $\partial_t\phi$ for the amplitudes of the plane wave. Its leading term is $(g_{aee}/m_e)(\nabla\phi)\cdot\vect{s}_e$, the spin-density row with a coefficient proportional to $|\vect{p}_\phi|$. Its second term contains $\op\psi^\dagger\boldsymbol{\sigma}\cdot\overleftrightarrow{\nabla}\op\psi$, whose matrix element between plane waves is $\ii\boldsymbol{\sigma}\cdot(\vect{k} + \vect{k}')$; it is the trace of the spin current, and the interaction becomes
\begin{equation}
  H_{aee} = \frac{g_{aee}}{m_e}\Bigl[(\nabla\phi)\cdot\vect{s}_e(\vect{p}_\phi) + (\partial_t\phi)\,\mathcal{J}^{\alpha\alpha}_e(\vect{p}_\phi)\Bigr] .
  \label{eq:8.1-alp}
\end{equation}
The trace couples $\vect{S}_j$ to the velocity of the same electron. With $|\partial_t\phi| = m_\phi|\phi|$ and $|\nabla\phi| = |\vect{p}_\phi||\phi|$, the second term exceeds the first by the electron velocity over the dark-matter velocity, as \cref{sec:1.4} found. When spin--orbit coupling is negligible the spin and the orbit of each electron separate, and the response of $\mathcal{J}^{\alpha\alpha}_e$ reduces to that of the current, which is why the rate again follows from the conductivity~\cite[Sect.~III\,B]{mitridate2021dark}.

The scalar couples through \cref{eq:1.4-scalar}. The operator $\op\psi^\dagger\overleftrightarrow{\nabla}^2\op\psi$ has the plane-wave matrix element $-(\vect{k} + \vect{k}')^2$, so $(1/8m_e^2)\op\psi^\dagger\overleftrightarrow{\nabla}^2\op\psi$ is $-\mathcal{T}_e/m_e$, and
\begin{keyresult}[Scalar absorption through the kinetic-energy density]
\begin{equation}
  H_{\phi ee} = -g_{\phi ee}\,\phi\,\Bigl[n_e(\vect{p}_\phi) - \frac{\mathcal{T}_e(\vect{p}_\phi)}{m_e}\Bigr] .
  \label{eq:8.1-scalar}
\end{equation}
The bracket is the non-relativistic limit of $\bar\psi\psi$: the number density times $1 - v^2/2$, the factor $(1 - v^2)^{1/2}$ by which a moving particle's scalar density falls below its number density. Matrix elements of $n_e(\vect{p}_\phi)$ between different states vanish as $\vect{p}_\phi \to 0$, and those of $\mathcal{T}_e$ do not, so the kinetic-energy term dominates absorption~\cite[Eq.~(18) and Sect.~III\,C]{mitridate2021dark}.
\end{keyresult}
\noindent The kinetic-energy row is the one that $\mathcal{O}_2$ needed, so a single new row per species covers both. The density term is suppressed twice: by orthogonality, since $\langle f|n_e(\vect{p}_\phi)|i\rangle \simeq \ii\vect{p}_\phi\cdot\langle f|\sum_j\vect{r}_j|i\rangle$, and by the screening of a density coupling in the medium~\cite[Sect.~III\,C]{mitridate2021dark}. The kinetic energy does not commute with the crystal Hamiltonian and has matrix elements between different states already at $\vect{p}_\phi = 0$. Its rate, unlike those of the dark photon and the axion-like particle, is not fixed by the optical conductivity, and \cref{sec:10.7} computes it.

Fermionic dark matter can be absorbed too. In the operators of Refs.~\cites[Eq.~(1)]{dror2020direct}[Eqs.~(1) and (2)]{dror2021absorption}, $\Lambda_F^{-2}(\bar{\tilde\nu}\Gamma\chi)(\bar\psi\Gamma'\psi)$ with a heavy scale $\Lambda_F$ and Dirac matrices $\Gamma$, $\Gamma'$, the dark-matter fermion $\chi$ is destroyed and a light fermion $\tilde\nu$ of mass $m_{\tilde\nu} \ll \mchi$ is emitted. The target bilinear $\bar\psi\Gamma'\psi$ reduces as for scattering: the vector to the density and current rows, the axial vector to the spin density and the trace of the spin current, the scalar to the bracket of \cref{eq:8.1-scalar}. What is new is the dark-matter side, which \cref{sec:8.2-absorption} builds into the rate, and the kinematics, which \cref{sec:8.1-kinematics} fixes.

\subsection{Energy deposited}\label{sec:8.1-kinematics}

In every process of this section the target receives the energy that the dark-matter side gives up. A particle of mass $\mchi$ and momentum $\vect{p}$ enters, and either nothing or a particle of momentum $\vect{p} - \vect{q}$ leaves, so that
\begin{equation}
  \omega = E_{\mathrm{in}}(\vect{p}) - E_{\mathrm{out}}(\vect{p} - \vect{q}) .
  \label{eq:8.1-omega}
\end{equation}
For elastic scattering, $E_{\mathrm{out}}$ is the kinetic energy of the same particle and \cref{eq:8.1-omega} is $\omega_{\vect{q}}$ of \cref{eq:1.3-omegaq} at $\Delta_\chi = 0$. For inelastic dark matter the outgoing state $\chi_2$ has mass $\mchi + \Delta_\chi$, and expanding its energy to first order in $(\vect{p} - \vect{q})^2$,
\begin{equation}
  \omega = \omega_{\vect{q}} + \frac{\Delta_\chi}{\mchi + \Delta_\chi}\,\frac{(\vect{p} - \vect{q})^2}{2\mchi} ,
  \label{eq:8.1-omega-inel}
\end{equation}
with $\omega_{\vect{q}}$ including the term $-\Delta_\chi$ (\cref{app:D.24}). The second term corrects the recoil energy by the fraction $\Delta_\chi/\mchi$, which \cref{eq:1.3-omegaq} drops; endothermic and exothermic scattering are then exactly the cases of \cref{sec:1.3}. A light outgoing fermion changes the picture. For $m_{\tilde\nu} \to 0$ and $v \to 0$, $E_{\mathrm{in}} = \mchi$ and $E_{\mathrm{out}} = |\vect{q}|$, so
\begin{equation}
  \omega = \mchi - q ,
  \label{eq:8.1-omega-ferm}
\end{equation}
and nearly all of the rest energy goes to $\tilde\nu$: the target is probed along a line in the $(q, \omega)$ plane at $q \approx \mchi$, far above the $q \le 2\mchi v$ of scattering. On a free nucleus of mass $M$ the recoil is then $\omega \approx \mchi^2/2M$~\cite[Eq.~(2)]{dror2020direct}, larger than an elastic recoil by $1/v^2$. Bosonic absorption is the limit in which nothing leaves, with $\omega = m_\phi$ at $\vect{q} = \vect{p}_\phi$. Scattering fills a region of the plane, fermionic absorption a line and bosonic absorption a point, and \cref{sec:8.2} shows that one rate formula, with the rows of \cref{eq:8.1-rows} and the tensor of \cref{eq:8.1-D}, covers all three.
